\documentclass[11pt,a4paper]{article}

% Langue et polices (compiler avec XeLaTeX)
\usepackage{fontspec}
\usepackage[french]{babel}

% Mathématiques, figures et liens
\usepackage{amsmath, amssymb}
\usepackage{graphicx}
\usepackage[hidelinks]{hyperref}

\title{Titre de l’article}
\author{Prénom Nom}
\date{\today}

\begin{document}

\maketitle

\begin{abstract}
  Bref résumé du travail : le problème, l’approche et les principaux résultats.
\end{abstract}

\section{Introduction}
\label{sec:introduction}

Ceci est un article d’exemple. Les accents, la cédille et les ligatures (é, è, ç, œ)
fonctionnent directement grâce à XeLaTeX. Pour plus de détails, voir la
section~\ref{sec:developpement} et la référence~\cite{knuth1984}.

\section{Développement}
\label{sec:developpement}

Une équation numérotée :
\begin{equation}
  \label{eq:euler}
  e^{i\pi} + 1 = 0
\end{equation}

et un système aligné :
\begin{align}
  \nabla \cdot \vec{E} &= \frac{\rho}{\varepsilon_0}, \\
  \nabla \cdot \vec{B} &= 0.
\end{align}

L’équation~\eqref{eq:euler} relie cinq constantes fondamentales.

\section{Conclusions}

Résumez ici les principales conclusions.

\bibliographystyle{plain}
\bibliography{references}

\end{document}
